\documentclass[../../main.tex]{subfiles}
\begin{document}

\begin{flushright}
  \footnotesize\textit{【自動文字起こし・要確認】}
\end{flushright}

{\bf [解]}
点$X$に対し，$\overrightarrow{OX}=\vec x$とおく．正四面体$OABC$の一辺を$1$として良く，
\begin{align}
  \vec p=\alpha\vec a，\quad \vec q=\beta\vec b，\quad \vec r=\gamma\vec c\quad(0<\alpha,\beta,\gamma<1)
\end{align}
とおく．$|PQ|=|QR|=|RP|$から
\begin{align}
  |\alpha\vec a-\beta\vec b|^2=|\beta\vec b-\gamma\vec c|^2=|\gamma\vec c-\alpha\vec a|^2
\end{align}

$|\vec a|=|\vec b|=|\vec c|=1$，$\vec a\cdot\vec b=\vec b\cdot\vec c=\vec c\cdot\vec a=\dfrac{1}{2}$だから，$k>0$として
\begin{align}
  \alpha^2-\alpha\beta+\beta^2&=k\quad\cdots①\\
  \beta^2-\beta\gamma+\gamma^2&=k\quad\cdots②\\
  \gamma^2-\gamma\alpha+\alpha^2&=k\quad\cdots③
\end{align}

①，②から
\begin{align}
  \alpha^2-\alpha\beta&=\gamma^2-\beta\gamma\\
  (\alpha^2-\gamma^2)-\beta(\alpha-\gamma)&=0\\
  (\alpha+\gamma-\beta)(\alpha-\gamma)&=0\quad\cdots④
\end{align}

同様に
\begin{align}
  (\alpha+\beta-\gamma)(\alpha-\beta)&=0\quad\cdots⑤\\
  (-\alpha+\beta+\gamma)(\beta-\gamma)&=0\quad\cdots⑥
\end{align}

④から，$\alpha=\gamma$または$\beta=\alpha+\gamma$である．

\underline{$1^\circ$ $\alpha=\gamma$}

⑤，⑥から$\beta=\alpha=\gamma$であり，この時$AB\parallel PQ$，$BC\parallel QR$，$CA\parallel RP$が成り立つ（$\because$相似）．

\underline{$2^\circ$ $\beta=\alpha+\gamma$}

⑤に代入して，$(\alpha+\beta-\gamma)(\alpha-\beta)=(2\alpha)(-\gamma)=0$．$\alpha\ne0$だから$\gamma=0$となり，$0<\gamma<1$に矛盾する．

以上から示された．$\blacksquare$

\begin{figure}[htb]
  \centering
  \begin{tikzpicture}[x={(3.2cm,0cm)}, y={(-1.2cm,-1cm)}, z={(0cm,2.2cm)},
    >=stealth, every node/.style={font=\small}]
    % 一辺 1 の正四面体の座標を投影する．
    \coordinate (A) at (0,0,0);
    \coordinate (B) at (1,0,0);
    \coordinate (C) at (0.5,0.8660254,0);
    \coordinate (O) at (0.5,0.2886751,0.8164966);
    \coordinate (P) at ($(O)!0.55!(A)$);
    \coordinate (Q) at ($(O)!0.55!(B)$);
    \coordinate (R) at ($(O)!0.55!(C)$);
    \fill[blue!12, opacity=0.65] (A)--(B)--(C)--cycle;
    \draw[dashed, gray!60, line width=0.9pt] (A)--(B);
    \draw[blue!85!black, line width=1.25pt]
    (A)--(C)--(B) (O)--(A) (O)--(B) (O)--(C);
    \fill[teal!20, opacity=0.7] (P)--(Q)--(R)--cycle;
    \draw[teal!85!black, line width=1.35pt] (P)--(Q)--(R)--cycle;
    \foreach \point/\position in {O/above,A/left,B/right,C/below} {
      \node[circle, fill=blue!85!black, inner sep=1.25pt] at (\point) {};
      \node[\position=2pt] at (\point) {$\point$};
    }
    \foreach \point/\position in {P/left,Q/right,R/left} {
      \node[circle, fill=red!85!black, inner sep=1.4pt] at (\point) {};
      \node[\position=2pt, red!85!black] at (\point) {$\point$};
    }
  \end{tikzpicture}
  \caption{正四面体$OABC$の辺上の点$P,Q,R$と平行な三角形の面}
\end{figure}

\newpage

\end{document}
